\documentclass[11pt]{article}

% ---- Packages ----
\usepackage[margin=1in]{geometry}
\usepackage{amsmath}
\usepackage{amssymb}
\usepackage{booktabs}
\usepackage{enumitem}
\usepackage{titlesec}
\usepackage[hidelinks]{hyperref}
\usepackage{xcolor}

% ---- Light formatting ----
\definecolor{moss}{HTML}{758C58}
\titleformat{\section}{\large\bfseries\color{moss}}{\thesection}{0.6em}{}
\titleformat{\subsection}{\normalsize\bfseries}{\thesubsection}{0.6em}{}
\setlist{itemsep=2pt, topsep=4pt}
\setlength{\parskip}{0.4em}
\setlength{\parindent}{0pt}

% ---- Title block ----
\title{\vspace{-1.5em}\textbf{Revealed Preferences of Large Language Models}\\[0.3em]
\large A Structural Estimation Approach to Model Behavior}
\author{J. Blakeslee \\ \small Pardee RAND Graduate School}
\date{\small Summer OJT Practicum (Course 1012) \\ \small Faculty: Todd Richmond \\ \small June 29 -- September 4, 2026}

\begin{document}
\maketitle

\hrule
\vspace{0.5em}

\section*{Abstract}
\noindent
This project treats a frontier large language model as an economic agent and recovers its
preference primitives from its choices. I elicit thousands of incentivized decisions from the
model under controlled experimental variation, then fit a structural discrete-choice model to
the resulting data. The warm-up estimates a single risk-aversion parameter from lottery choices.
The main study estimates loss aversion and a reference point from systematically framed gains and
losses, testing whether the model's behavior obeys reference-dependent preferences in the sense of
Kahneman and Tversky and Kőszegi and Rabin. The deliverable is a reproducible pipeline,
a clean dataset of model choices, and a short working paper reporting estimated parameters with
confidence intervals. The central finding will be a quantitative answer to one question: does a
frontier model carry stable, human-like behavioral preferences, and do those preferences survive
reframing?

\vspace{0.5em}
\hrule
\vspace{0.8em}

\section{The Research Question}

The project answers one question and a tightly related follow-up:

\begin{enumerate}[label=\textbf{Q\arabic*.}, leftmargin=*]
  \item \textbf{Do frontier language models exhibit stable, structured economic preferences?}
  Concretely: when a model makes risky choices, can its behavior be summarized by interpretable
  primitives (a risk-aversion coefficient, a loss-aversion parameter, a reference point), the way
  a human subject's choices can?
  \item \textbf{Are those preferences consistent under reframing?} The same underlying choice, framed
  once as a gain and once as a loss, should produce the same decision from a rational agent. Humans
  famously violate this. Does the model?
\end{enumerate}

The answer is a small set of estimated numbers with standard errors and a clear verdict: the model
is risk-averse with coefficient $r$; its loss-aversion parameter $\lambda$ is (or is not)
statistically distinguishable from the rational benchmark of one; its choices do (or do not) flip
when the frame moves. Each of those is a real, reportable result.

\section{Why This Question Matters}

\subsection{Models are becoming economic agents}
Language models increasingly make allocative decisions in deployment: pricing, resource allocation,
negotiation, portfolio and risk judgments inside agentic systems. The behavior of an agent that
acts on the world is a governance object. Before we can reason about an automated agent's market or
policy effects, we need to know what it actually prefers when it chooses. Revealed-preference
estimation is the discipline economics built for exactly this purpose, and it has barely been
pointed at models.

\subsection{Behavioral consistency is an alignment property}
A model whose choices reverse when a problem is reframed is, in a precise sense, manipulable. The
same susceptibility that makes a human buy the extended warranty makes a model exploitable by an
adversarial prompt or an interface designed to steer it. Measuring whether a model's preferences are
frame-invariant is therefore a safety and evaluations result, stated in the rigorous language of
choice theory rather than as an anecdote. ``Does this model's risk attitude survive reframing'' is a
real robustness question that current evals largely do not answer.

\subsection{It tells us what alignment training does to behavior}
If a base model is roughly frame-invariant and a fine-tuned model exhibits human-like loss aversion,
that is evidence that reinforcement learning from human feedback transferred human cognitive biases
into the model along with human values. If the reverse holds, that is equally informative. Either way,
the estimate speaks to a question alignment teams care about and rarely quantify: what kind of agent
does post-training actually produce?

\subsection{The contribution is structural, which buys counterfactuals}
Showing that a model behaves oddly in one experiment is a description. Recovering the primitives
behind its behavior lets me predict its choices in environments I never ran. That is the difference
between ``the model gambled here'' and ``the model has a CRRA coefficient of $r$, so it will do the
following in any lottery you hand it.'' The structural step is what makes the result generalize, and
it is the part most adjacent work omits.

\subsection{Positioning against the existing literature}
A growing body of work asks whether models display human biases (Horton's ``homo silicus,''
Binz and Schulz's cognitive-psychology probes, Mei and coauthors' behavioral Turing test). That
literature largely documents biases qualitatively. The gap, and this project's niche, is the
\emph{structural estimation} of reference-dependent preferences with a likelihood, confidence
intervals, and an explicit test against the rational benchmark, combined with the question of whether
post-training instilled the bias. The aim is a measured parameter, not another demonstration that
models can be nudged.

\section{Methodology}

The project runs in two phases. The first is a deliberately easy warm-up that de-risks the entire
pipeline. The second is the real study and reuses the same machinery with one added primitive.

\subsection{Phase 1 (warm-up): risk preferences from lottery menus}
Following the multiple-price-list tradition (Holt and Laury; Andersen, Harrison, Lau, and Rutström),
I present the model with paired gambles across a grid of stakes and probabilities, for example a sure
\$50 against a 50\% chance of \$120. Each trial is an independent call with no memory of prior choices.
I model utility as constant relative risk aversion,
\[
  u(x) = \frac{x^{1-r}}{1-r},
\]
and assume a logit (Fechner) choice error, so the probability of choosing lottery $B$ over option $A$ is
\[
  \Pr(B) = \frac{1}{1 + \exp\!\big(-\,(EU_B - EU_A)/\mu\big)},
\]
where $EU$ is expected utility and $\mu$ is a noise parameter. I estimate $r$ and $\mu$ by maximum
likelihood. This is the canonical structural exercise; its answer can be checked against any
first-year problem set, which is precisely why it is the right place to debug elicitation before the
stakes rise.

\subsection{Phase 2 (main study): reference dependence and loss aversion}
The same lottery is now presented twice: once framed as a gain from a low reference point and once as
a loss from a high reference point (``you start with \$120; option $B$ risks losing all of it''). The
underlying payoffs are identical; only the frame moves. I extend the utility specification to carry a
reference point $\rho$ and a loss-aversion parameter $\lambda$, in the spirit of Tversky and Kahneman
and Kőszegi and Rabin:
\[
  v(x \mid \rho) =
  \begin{cases}
    (x - \rho)^{\alpha}, & x \ge \rho \\[4pt]
    -\lambda\,(\rho - x)^{\beta}, & x < \rho.
  \end{cases}
\]
I estimate $\lambda$, the reference point's sensitivity to the frame, and the curvature terms, again
by maximum likelihood on the binary choices. The headline tests are sharp hypotheses:
\begin{itemize}
  \item $H_0: \lambda = 1$ (no loss aversion; the rational benchmark).
  \item $H_0:$ the estimated reference point does not move with the frame (frame invariance).
\end{itemize}
Rejecting the first says the model is loss-averse. Rejecting the second says its choices are
manipulable by framing. Failing to reject either is itself a clean and reportable result.

\subsection{Elicitation hygiene: where the new skill lives}
Valid estimation depends on a clean dataset, and producing one from a non-human subject is the genuinely
new craft this project teaches. Four controls matter:
\begin{itemize}
  \item \textbf{Forced parseable output.} The model is constrained to emit a single option label, via
  structured output or a first-token rule, and every response is validated; malformed responses are
  logged and re-run, with the discard rate tracked as a diagnostic.
  \item \textbf{Randomized labels and order.} The mapping from option letters to gambles, and the order
  in which options appear, are randomized per trial and recorded, so a positional or label bias cannot
  masquerade as a preference.
  \item \textbf{Stochasticity by design.} Running at a moderate temperature makes choices genuinely
  stochastic, which gives the logit error term its meaning and supplies within-condition variance for
  estimation.
  \item \textbf{Independence.} Each trial is a fresh context with no conversation history, so the draws
  are independent and the static choice model is valid.
\end{itemize}
A full Phase 2 dataset is roughly 40 gambles $\times$ 2 frames $\times$ 25 repetitions, on the order of
2{,}000 calls, which costs a few dollars and runs in under an hour. The generous repetition count is what
makes the confidence intervals tight.

\section{What I Will Learn}

This project is chosen as a skill-building artifact. It sits one deliberate notch past my current
toolkit on two axes at once.

\begin{itemize}
  \item \textbf{Model elicitation and experimental design on AI systems.} Designing controlled
  experiments for a model rather than a human: API orchestration, output forcing, randomization
  hygiene, the role of temperature, and independence. This is the foreign skill, and it is the one
  that makes a frontier lab read the work as ``this person can design and run model evaluations.''
  \item \textbf{Structural estimation of preference primitives.} My empirical training is reduced-form
  and causal (difference-in-differences, synthetic control, instrumental variables). Estimating a
  behavioral model's primitives from choice data, then doing counterfactuals the data never saw, is
  the structural step beyond that. It is the credential that is hard to fake and that compounds with
  everything else I do.
  \item \textbf{The bridge between behavioral economics and AI evaluation.} Learning to state an
  alignment-relevant property (frame invariance, manipulability) as a hypothesis test on estimated
  parameters, so that a behavioral question becomes a measured result.
  \item \textbf{Reproducible research engineering.} A clean, public pipeline from elicitation to dataset
  to estimation to writeup, the kind of artifact that is legible and re-runnable by a reviewer.
\end{itemize}

\section{Timeline}

\begin{center}
\begin{tabular}{@{}ll@{}}
\toprule
\textbf{Weeks} & \textbf{Milestone} \\
\midrule
1--2 & Literature scan; build the elicitation harness; first lottery dataset out of a model. \\
3    & Estimate Phase 1 risk aversion; warm-up result and full pipeline locked. \\
4--5 & Design Phase 2: framing templates, reference-point manipulation, hypothesis structure. \\
6    & Run Phase 2 elicitation; assemble the reference-dependence dataset. \\
7--8 & Estimate loss aversion and reference point; run the two sharp hypothesis tests. \\
9    & Robustness (temperature, label randomization, a second model family if time); draft writeup. \\
10   & Final briefing to the cohort; publish a clean, reproducible notebook and short paper. \\
\bottomrule
\end{tabular}
\end{center}

A presentable result exists by week three regardless of how Phase 2 unfolds, so the finished
deliverable is secured early and the remaining weeks add ambition rather than risk.

\section{Deliverables}

\begin{itemize}
  \item A reproducible code repository: elicitation harness plus estimation, runnable end to end.
  \item A clean, documented dataset of the model's choices under every condition.
  \item A short working paper reporting the estimated parameters, confidence intervals, and the verdict
  on consistency under reframing.
  \item The final sprint briefing, with a public-facing explainer of the headline finding.
\end{itemize}

\section{Fit with the Sprint}

The practicum asks for applications work that informs implications work. This project is exactly that
loop in miniature: the application is a working model-evaluation pipeline I build and run; the
implication is a governance-relevant claim about whether frontier agents carry stable, manipulable
preferences. The technical build produces the evidence, and the evidence produces the argument. The
result is legible across the audiences the sprint serves: a lab evaluations or economics team reads it
as rigorous model elicitation; a policy reader reads it as a concrete statement about how an automated
economic agent behaves under pressure.

\vspace{1em}
\hrule
\vspace{0.5em}

\begin{thebibliography}{9}
\small

\bibitem{andersen2008}
Andersen, S., Harrison, G. W., Lau, M. I., and Rutström, E. E. (2008).
Eliciting Risk and Time Preferences. \textit{Econometrica}, 76(3), 583--618.

\bibitem{binz2023}
Binz, M., and Schulz, E. (2023).
Using Cognitive Psychology to Understand GPT-3. \textit{Proceedings of the National Academy of Sciences}, 120(6).

\bibitem{holt2002}
Holt, C. A., and Laury, S. K. (2002).
Risk Aversion and Incentive Effects. \textit{American Economic Review}, 92(5), 1644--1655.

\bibitem{horton2023}
Horton, J. J. (2023).
Large Language Models as Simulated Economic Agents: What Can We Learn from Homo Silicus?
\textit{NBER Working Paper} 31122.

\bibitem{kahneman1979}
Kahneman, D., and Tversky, A. (1979).
Prospect Theory: An Analysis of Decision under Risk. \textit{Econometrica}, 47(2), 263--291.

\bibitem{koszegi2006}
Kőszegi, B., and Rabin, M. (2006).
A Model of Reference-Dependent Preferences. \textit{Quarterly Journal of Economics}, 121(4), 1133--1165.

\bibitem{mei2024}
Mei, Q., Xie, Y., Yuan, W., and Jackson, M. O. (2024).
A Turing Test of Whether AI Chatbots Are Behaviorally Similar to Humans.
\textit{Proceedings of the National Academy of Sciences}, 121(9).

\bibitem{train2009}
Train, K. E. (2009).
\textit{Discrete Choice Methods with Simulation} (2nd ed.). Cambridge University Press.

\bibitem{tversky1991}
Tversky, A., and Kahneman, D. (1991).
Loss Aversion in Riskless Choice: A Reference-Dependent Model.
\textit{Quarterly Journal of Economics}, 106(4), 1039--1061.

\end{thebibliography}

\end{document}
